% !TeX root = part02.tex

\documentclass[../final.tex]{subfiles}

\begin{document}

% ============================================================
% Семинар 2
% ============================================================

\practice{2}{11.09.2026}{Вырожденный ферми-газ и холодные звёзды}


% ============================================================
% Задача 1
% ============================================================

\section{Задача 1. Оценка размера многоэлектронного атома}

\textbf{Условие.}
В атомах с большим числом $Z$ электронов большая часть из них
движется вблизи ядра в объёме с характерным размером $R$.
В модели Томаса--Ферми эти электроны рассматриваются как
вырожденный газ, температура которого равна нулю.
Пренебрегая кулоновским взаимодействием электронов друг с другом,
оценить величину $R$.


% ------------------------------------------------------------
% Схема атома
% ------------------------------------------------------------

\begin{rbox}[left=18pt,right=18pt,top=14pt,bottom=10pt]{[]}
    \centering

    \begin{tikzpicture}[
        >=Latex,
        every node/.style={text=rtext}
    ]

        % Область электронной оболочки
        \draw[
            draw=rc3,
            line width=1.0pt
        ]
        (0,0) circle (2.15);

        % Ядро
        \fill[rc3] (0,0) circle (0.18);

        % Электроны
        \foreach \x/\y in {
            -1.55/0.55,
            -1.30/-0.75,
            -0.95/1.35,
            -0.55/-1.55,
            -0.25/0.90,
             0.30/-1.20,
             0.55/1.55,
             0.95/0.65,
             1.25/-0.85,
             1.60/0.20,
            -1.75/-0.10,
             1.05/1.30
        }{
            \fill[rtext] (\x,\y) circle (1.5pt);
        }

        % Радиус
        \draw[
            ->,
            draw=rtext,
            line width=0.9pt
        ]
        (0,0) -- (1.90,0);

        \node[
            above
        ]
        at (0.95,0)
        {$R$};

        \node[
            below
        ]
        at (0,-0.22)
        {\small $+Ze$};

    \end{tikzpicture}

    \captionsetup{
        font=footnotesize,
        labelfont=bf,
        skip=5pt
    }

    \captionof{figure}{
        Оценочная модель многоэлектронного атома:
        ядро с зарядом $+Ze$ и вырожденный электронный газ,
        локализованный в области характерного радиуса $R$.
    }
    \label{fig:statphys_tf_atom}
\end{rbox}


% ------------------------------------------------------------
% Решение
% ------------------------------------------------------------

\subsection*{Решение}

При $T=0$ электронный газ является полностью вырожденным.
Полную энергию атома оценим как сумму кинетической энергии
электронов и их потенциальной энергии в поле ядра:
%
\begin{equation*}
    E
    =
    E_{\text{кин}}
    +
    E_{\text{пот}}.
\end{equation*}


% ------------------------------------------------------------
% Кинетическая энергия
% ------------------------------------------------------------

\subsection*{Кинетическая энергия электронов}

Для нерелятивистского вырожденного электронного газа
%
\begin{equation*}
    \varepsilon_F
    =
    \frac{p_F^2}{2m_e}.
\end{equation*}

По порядку величины импульс Ферми равен
%
\begin{equation*}
    p_F
    \sim
    \hbar n^{1/3}.
\end{equation*}

В атоме имеется $Z$ электронов, заключённых в области объёма
порядка
%
\begin{equation*}
    V
    \sim
    R^3.
\end{equation*}

Поэтому концентрация электронов имеет порядок
%
\begin{equation*}
    n
    \sim
    \frac{Z}{R^3}.
\end{equation*}

Следовательно,
%
\begin{equation*}
    p_F
    \sim
    \hbar
    \left(
        \frac{Z}{R^3}
    \right)^{1/3}
    =
    \frac{\hbar Z^{1/3}}{R}.
\end{equation*}

Энергия Ферми тогда равна по порядку величины
%
\begin{equation*}
    \varepsilon_F
    \sim
    \frac{\hbar^2 Z^{2/3}}{2m_eR^2}.
\end{equation*}

Поскольку электронов всего $Z$, полная кинетическая энергия
имеет порядок
%
\begin{equation*}
    E_{\text{кин}}
    \sim
    Z\varepsilon_F
    \sim
    \frac{\hbar^2Z^{5/3}}{2m_eR^2}.
\end{equation*}


% ------------------------------------------------------------
% Потенциальная энергия
% ------------------------------------------------------------

\subsection*{Потенциальная энергия}

Потенциальная энергия одного электрона в поле ядра с зарядом $+Ze$
при характерном расстоянии $R$ имеет порядок
%
\begin{equation*}
    U
    \sim
    -\frac{Ze^2}{R}.
\end{equation*}

Для всех $Z$ электронов
%
\begin{equation*}
    E_{\text{пот}}
    \sim
    -\frac{Z^2e^2}{R}.
\end{equation*}

Здесь, как и в рукописном решении, используется гауссова система
единиц, поэтому множитель $1/(4\pi\varepsilon_0)$ явно не пишется.

Таким образом,
%
\begin{equation*}
    E(R)
    \sim
    \frac{\hbar^2Z^{5/3}}{2m_eR^2}
    -
    \frac{Z^2e^2}{R}.
\end{equation*}


% ------------------------------------------------------------
% Минимум энергии
% ------------------------------------------------------------

\subsection*{Равновесный размер атома}

Равновесному размеру соответствует минимум полной энергии:
%
\begin{equation*}
    \frac{dE}{dR}
    =
    0.
\end{equation*}

Дифференцируя,
%
\begin{equation*}
    -\frac{\hbar^2Z^{5/3}}{m_eR^3}
    +
    \frac{Z^2e^2}{R^2}
    =
    0.
\end{equation*}

Умножая на $R^3$, получаем
%
\begin{equation*}
    -\frac{\hbar^2Z^{5/3}}{m_e}
    +
    Z^2e^2R
    =
    0.
\end{equation*}

Следовательно,
%
\begin{equation*}
    Z^2e^2R
    =
    \frac{\hbar^2Z^{5/3}}{m_e},
\end{equation*}
%
откуда
%
\begin{equation*}
    R
    \sim
    \frac{\hbar^2}{m_ee^2}
    Z^{-1/3}.
\end{equation*}

Введём боровский радиус
%
\begin{equation*}
    a_0
    =
    \frac{\hbar^2}{m_ee^2}.
\end{equation*}

Тогда окончательно
%
\begin{equation*}
    \boxed{
        R
        \sim
        \frac{a_0}{Z^{1/3}}
    }.
\end{equation*}

Таким образом, при увеличении заряда ядра характерный размер
внутренней электронной области уменьшается как
%
\begin{equation*}
    R
    \propto
    Z^{-1/3}.
\end{equation*}


% ============================================================
% Задача 2
% ============================================================

\newpage
\section{Задача 2. Равновесный радиус белого карлика}

\textbf{Условие.}
Когда в обычной звезде выгорает заметная часть водорода,
термоядерная реакция прекращается и звезда, излучая, будет остывать
и сжиматься. При сжатии гравитационная энергия уменьшается,
однако энергия вырожденного электронного газа возрастает.

В результате холодная звезда, называемая белым карликом,
достигает некоторого равновесного радиуса.
Оценить радиус белого карлика из условия минимума полной энергии.

Использовать
%
\begin{equation*}
    M_\star
    =
    2\times10^{33}\,\text{г},
    \qquad
    G
    =
    6.6\times10^{-8}\,
    \frac{\text{дин}\cdot\text{см}^2}{\text{г}^2}.
\end{equation*}


% ------------------------------------------------------------
% Схема белого карлика
% ------------------------------------------------------------

\begin{rbox}[left=18pt,right=18pt,top=14pt,bottom=10pt]{[]}
    \centering

    \begin{tikzpicture}[
        >=Latex,
        every node/.style={text=rtext}
    ]

        % Звезда
        \draw[
            draw=rc3,
            line width=1.0pt
        ]
        (0,0) circle (1.55);

        \fill[rc3] (0,0) circle (2.0pt);

        % Гравитация: стрелки внутрь
        \foreach \a in {0,90,180,270}{
            \draw[
                ->,
                draw=rtext,
                line width=0.8pt
            ]
            ({2.35*cos(\a)},{2.35*sin(\a)})
            --
            ({1.62*cos(\a)},{1.62*sin(\a)});
        }

        % Давление вырождения: стрелки наружу
        \foreach \a in {45,135,225,315}{
            \draw[
                ->,
                draw=rc3,
                line width=0.9pt
            ]
            ({1.00*cos(\a)},{1.00*sin(\a)})
            --
            ({1.75*cos(\a)},{1.75*sin(\a)});
        }

        % Радиус
        \draw[
            ->,
            draw=rtext,
            line width=0.8pt
        ]
        (0,0) -- (1.42,0);

        \node[
            above
        ]
        at (0.72,0)
        {$R$};

        \node[
            right
        ]
        at (2.42,0)
        {\small гравитация};

        \node[
            right
        ]
        at (1.55,1.42)
        {\small давление вырождения};

    \end{tikzpicture}

    \captionsetup{
        font=footnotesize,
        labelfont=bf,
        skip=5pt
    }

    \captionof{figure}{
        Схематическое равновесие белого карлика:
        гравитационное сжатие компенсируется давлением
        вырожденного электронного газа.
    }
    \label{fig:statphys_white_dwarf}
\end{rbox}


% ------------------------------------------------------------
% Полная энергия
% ------------------------------------------------------------

\subsection*{Полная энергия звезды}

Будем считать звезду холодной, поэтому электронный газ практически
полностью вырожден.

Полная энергия имеет вид
%
\begin{equation*}
    E
    =
    E_{\text{грав}}
    +
    E_{\text{кин},e}
    +
    E_{\text{кин},я},
\end{equation*}
%
где $E_{\text{кин},e}$ --- кинетическая энергия электронов,
а $E_{\text{кин},я}$ --- кинетическая энергия ядер.


% ------------------------------------------------------------
% Гравитационная энергия
% ------------------------------------------------------------

\subsection*{Гравитационная энергия}

По порядку величины гравитационная энергия звезды массы $M$
и радиуса $R$ равна
%
\begin{equation*}
    E_{\text{грав}}
    \sim
    -\frac{GM^2}{R}.
\end{equation*}

Численный коэффициент порядка единицы здесь несущественен,
поскольку требуется оценка.


% ------------------------------------------------------------
% Число электронов
% ------------------------------------------------------------

\subsection*{Число электронов}

Основная часть массы звезды заключена в ядрах.
По порядку величины число электронов можно оценить как
%
\begin{equation*}
    N_e
    \sim
    \frac{M}{m_p}.
\end{equation*}

Более точный коэффициент зависит от химического состава белого карлика,
но для оценки он несущественен.


% ------------------------------------------------------------
% Почему вклад электронов главный
% ------------------------------------------------------------

\subsection*{Электроны и ядра}

Для нерелятивистского вырожденного газа энергия Ферми имеет вид
%
\begin{equation*}
    \varepsilon_F
    =
    \frac{p_F^2}{2m}.
\end{equation*}

При сравнимых концентрациях характерные импульсы Ферми электронов
и ядер имеют один порядок, но
%
\begin{equation*}
    m_e
    \ll
    m_{\text{я}}.
\end{equation*}

Следовательно,
%
\begin{equation*}
    \varepsilon_{F,e}
    \gg
    \varepsilon_{F,\text{я}},
\end{equation*}
%
и поэтому
%
\begin{equation*}
    E_{\text{кин},e}
    \gg
    E_{\text{кин},\text{я}}.
\end{equation*}

В дальнейшем вкладом ядер в кинетическую энергию пренебрегаем.


% ------------------------------------------------------------
% Кинетическая энергия электронов
% ------------------------------------------------------------

\subsection*{Кинетическая энергия вырожденных электронов}

Энергия Ферми электронов
%
\begin{equation*}
    \varepsilon_{F,e}
    \sim
    \frac{\hbar^2n_e^{2/3}}{2m_e}.
\end{equation*}

Полная кинетическая энергия имеет порядок
%
\begin{equation*}
    E_{\text{кин},e}
    \sim
    N_e\varepsilon_{F,e}.
\end{equation*}

Следовательно,
%
\begin{equation*}
    E_{\text{кин},e}
    \sim
    \frac{\hbar^2n_e^{2/3}}{2m_e}
    N_e.
\end{equation*}

Концентрация электронов равна по порядку величины
%
\begin{equation*}
    n_e
    \sim
    \frac{N_e}{R^3}
    \sim
    \frac{M}{m_pR^3}.
\end{equation*}

Поэтому
%
\begin{equation*}
    n_e^{2/3}
    \sim
    \frac{1}{R^2}
    \left(
        \frac{M}{m_p}
    \right)^{2/3}.
\end{equation*}

Используя
%
\begin{equation*}
    N_e
    \sim
    \frac{M}{m_p},
\end{equation*}
%
получаем
%
\begin{equation*}
    E_{\text{кин},e}
    \sim
    \frac{\hbar^2}{2m_eR^2}
    \left(
        \frac{M}{m_p}
    \right)^{5/3}.
\end{equation*}

Таким образом, полная энергия белого карлика оценивается как
%
\begin{equation*}
    E(R)
    \sim
    -\frac{GM^2}{R}
    +
    \frac{\hbar^2}{2m_eR^2}
    \left(
        \frac{M}{m_p}
    \right)^{5/3}.
\end{equation*}


% ------------------------------------------------------------
% Минимум полной энергии
% ------------------------------------------------------------

\subsection*{Равновесный радиус}

Равновесие соответствует минимуму полной энергии:
%
\begin{equation*}
    \frac{dE}{dR}
    =
    0.
\end{equation*}

Дифференцируя,
%
\begin{equation*}
    \frac{GM^2}{R^2}
    -
    \frac{\hbar^2}{m_eR^3}
    \left(
        \frac{M}{m_p}
    \right)^{5/3}
    =
    0.
\end{equation*}

Умножаем на $R^3$:
%
\begin{equation*}
    GM^2R
    =
    \frac{\hbar^2}{m_e}
    \left(
        \frac{M}{m_p}
    \right)^{5/3}.
\end{equation*}

Отсюда
%
\begin{equation*}
    R
    \sim
    \frac{\hbar^2}{Gm_e}
    \frac{(M/m_p)^{5/3}}{M^2}.
\end{equation*}

Поскольку
%
\begin{equation*}
    \frac{M^{5/3}}{M^2}
    =
    M^{-1/3},
\end{equation*}
%
получаем
%
\begin{equation*}
    \boxed{
        R
        \sim
        \frac{\hbar^2}
        {Gm_em_p^{5/3}}
        M^{-1/3}
    }.
\end{equation*}

Таким образом,
%
\begin{equation*}
    \boxed{
        R
        \propto
        M^{-1/3}
    }.
\end{equation*}

Чем больше масса белого карлика, тем меньше его равновесный радиус.


% ------------------------------------------------------------
% Численная оценка
% ------------------------------------------------------------

\subsection*{Численная оценка}

Используем
%
\begin{equation*}
    \hbar
    \simeq
    1.05\times10^{-27}\,
    \text{эрг}\cdot\text{с},
\end{equation*}
%
\begin{equation*}
    m_e
    \simeq
    9.1\times10^{-28}\,\text{г},
    \qquad
    m_p
    \simeq
    1.67\times10^{-24}\,\text{г},
\end{equation*}
%
\begin{equation*}
    M
    \simeq
    2\times10^{33}\,\text{г}.
\end{equation*}

Тогда по порядку величины
%
\begin{equation*}
    R
    \sim
    6\times10^8\,\text{см}.
\end{equation*}

Или
%
\begin{equation*}
    \boxed{
        R
        \sim
        6\times10^3\,\text{км}
    }.
\end{equation*}

Полученный размер имеет порядок радиуса Земли, что характерно
для белых карликов.


% ============================================================
% Задача 3
% ============================================================

\newpage
\section{Задача 3. Ультрарелятивистский вырожденный газ}

\textbf{Условие.}
Если звезда на начальной стадии имела массу, заметно большую массы
Солнца, или масса белого карлика может увеличиваться за счёт захвата
окружающего вещества, то газ электронов может стать
ультрарелятивистским.

Оценить предельное значение массы белого карлика.

Найти давление ультрарелятивистского газа.

Найти соотношение между концентрациями электронов и нейтронов,
считая их ультрарелятивистскими.


% ------------------------------------------------------------
% Релятивистская кинематика
% ------------------------------------------------------------

\subsection*{Релятивистская энергия частицы}

Для свободной релятивистской частицы
%
\begin{equation*}
    \varepsilon^2
    =
    p^2c^2
    +
    m^2c^4.
\end{equation*}

Или
%
\begin{equation*}
    \frac{\varepsilon^2}{c^2}
    =
    p^2
    +
    m^2c^2.
\end{equation*}

Лагранжиан свободной частицы имеет вид
%
\begin{equation*}
    L
    =
    -mc^2
    \sqrt{
        1-\frac{v^2}{c^2}
    }.
\end{equation*}

Соответствующее действие можно записать в инвариантной форме:
%
\begin{equation*}
    S
    =
    -mc
    \int ds,
\end{equation*}
%
где
%
\begin{equation*}
    ds^2
    =
    c^2dt^2
    -
    d\mathbf r^2.
\end{equation*}

Импульс частицы
%
\begin{equation*}
    \mathbf p
    =
    \frac{\partial L}{\partial\mathbf v}
    =
    \frac{m\mathbf v}
    {\sqrt{1-\beta^2}},
    \qquad
    \beta
    =
    \frac{v}{c}.
\end{equation*}

Энергия определяется преобразованием Лежандра:
%
\begin{equation*}
    \varepsilon
    =
    \mathbf v\mathbf p
    -
    L.
\end{equation*}

Подставляя $\mathbf p$ и $L$, получаем
%
\begin{align*}
    \varepsilon
    &=
    \frac{mv^2}{\sqrt{1-\beta^2}}
    +
    mc^2\sqrt{1-\beta^2}
    \\[3pt]
    &=
    \frac{
        mv^2+mc^2(1-\beta^2)
    }{
        \sqrt{1-\beta^2}
    }
    =
    \frac{mc^2}{\sqrt{1-\beta^2}}.
\end{align*}

Следовательно,
%
\begin{equation*}
    \boxed{
        \varepsilon
        =
        \frac{mc^2}{\sqrt{1-\beta^2}}
    }.
\end{equation*}


% ------------------------------------------------------------
% Импульс Ферми
% ------------------------------------------------------------

\subsection*{Импульс Ферми}

При $T=0$ электроны заполняют все состояния с импульсами
%
\begin{equation*}
    0
    \leq
    p
    \leq
    p_F.
\end{equation*}

С учётом двух проекций спина
%
\begin{equation*}
    N_e
    =
    2V
    \int_0^{p_F}
    \frac{d^3p}{(2\pi\hbar)^3}.
\end{equation*}

В сферических координатах
%
\begin{equation*}
    d^3p
    =
    4\pi p^2\,dp.
\end{equation*}

Поэтому
%
\begin{equation*}
    N_e
    =
    \frac{8\pi V}{(2\pi\hbar)^3}
    \int_0^{p_F}
    p^2\,dp.
\end{equation*}

Так как
%
\begin{equation*}
    \int_0^{p_F}
    p^2\,dp
    =
    \frac{p_F^3}{3},
\end{equation*}
%
получаем
%
\begin{equation*}
    N_e
    =
    \frac{Vp_F^3}{3\pi^2\hbar^3}.
\end{equation*}

Следовательно,
%
\begin{equation*}
    n_e
    =
    \frac{N_e}{V}
    =
    \frac{p_F^3}{3\pi^2\hbar^3},
\end{equation*}
%
откуда
%
\begin{equation*}
    \boxed{
        p_F
        =
        \hbar
        \left(
            3\pi^2n_e
        \right)^{1/3}
    }.
\end{equation*}

Для оценок будем использовать более простую запись
%
\begin{equation*}
    p_F
    \sim
    \hbar n_e^{1/3}.
\end{equation*}


% ------------------------------------------------------------
% Ультрарелятивистский предел
% ------------------------------------------------------------

\subsection*{Ультрарелятивистский предел}

Если
%
\begin{equation*}
    p_F
    \gg
    m_ec,
\end{equation*}
%
то массой электрона в законе дисперсии можно пренебречь:
%
\begin{equation*}
    \varepsilon
    \simeq
    pc.
\end{equation*}

В частности,
%
\begin{equation*}
    \varepsilon_F
    \simeq
    p_Fc.
\end{equation*}

Полная кинетическая энергия электронов по порядку величины равна
%
\begin{equation*}
    E_{\text{кин}}
    \sim
    N_e\varepsilon_F
    \sim
    N_ep_Fc.
\end{equation*}

Так как
%
\begin{equation*}
    p_F
    \sim
    \hbar
    \left(
        \frac{N_e}{R^3}
    \right)^{1/3}
    =
    \frac{\hbar N_e^{1/3}}{R},
\end{equation*}
%
то
%
\begin{equation*}
    E_{\text{кин}}
    \sim
    \frac{\hbar c\,N_e^{4/3}}{R}.
\end{equation*}

Для белого карлика
%
\begin{equation*}
    N_e
    \sim
    \frac{M}{m_p},
\end{equation*}
%
поэтому
%
\begin{equation*}
    E_{\text{кин}}
    \sim
    \frac{\hbar c}{R}
    \left(
        \frac{M}{m_p}
    \right)^{4/3}.
\end{equation*}


% ------------------------------------------------------------
% Предельная масса
% ------------------------------------------------------------

\subsection*{Предельная масса белого карлика}

Полная энергия звезды теперь имеет вид
%
\begin{equation*}
    E(R)
    \sim
    -\frac{GM^2}{R}
    +
    \frac{\hbar c}{R}
    \left(
        \frac{M}{m_p}
    \right)^{4/3}.
\end{equation*}

Вынесем $1/R$:
%
\begin{equation*}
    E(R)
    \sim
    \frac{1}{R}
    \left[
        \hbar c
        \left(
            \frac{M}{m_p}
        \right)^{4/3}
        -
        GM^2
    \right].
\end{equation*}

В отличие от нерелятивистского случая, обе части энергии теперь
зависят от радиуса одинаково:
%
\begin{equation*}
    E_{\text{кин}}
    \propto
    \frac{1}{R},
    \qquad
    E_{\text{грав}}
    \propto
    -\frac{1}{R}.
\end{equation*}

Поэтому устойчивый минимум при конечном $R$ исчезает.

Чтобы электронное давление ещё могло противостоять гравитации,
необходимо
%
\begin{equation*}
    \hbar c
    \left(
        \frac{M}{m_p}
    \right)^{4/3}
    \gtrsim
    GM^2.
\end{equation*}

Переносим степени массы:
%
\begin{equation*}
    \frac{\hbar c}{m_p^{4/3}}
    \gtrsim
    GM^{2/3}.
\end{equation*}

Следовательно,
%
\begin{equation*}
    M^{2/3}
    \lesssim
    \frac{\hbar c}
    {Gm_p^{4/3}}.
\end{equation*}

Возводя обе части в степень $3/2$, получаем
%
\begin{equation*}
    \boxed{
        M
        \lesssim
        \left(
            \frac{\hbar c}{G}
        \right)^{3/2}
        \frac{1}{m_p^2}
    }.
\end{equation*}

Таким образом, характерная предельная масса равна
%
\begin{equation*}
    \boxed{
        M_{\text{кр}}
        \sim
        \frac{
            (\hbar c/G)^{3/2}
        }{
            m_p^2
        }
    }.
\end{equation*}

Численно
%
\begin{equation*}
    M_{\text{кр}}
    \sim
    4\times10^{33}\,\text{г},
\end{equation*}
%
то есть величина имеет порядок нескольких масс Солнца:
%
\begin{equation*}
    \boxed{
        M_{\text{кр}}
        \sim
        M_\odot
    }.
\end{equation*}

Это оценка по порядку величины.
Точное значение предела Чандрасекара зависит от состава вещества,
профиля плотности и численных коэффициентов, отброшенных в данной оценке.


% ------------------------------------------------------------
% Иллюстрация устойчивости
% ------------------------------------------------------------

\begin{rbox}[left=18pt,right=18pt,top=14pt,bottom=10pt]{[]}
    \centering

    \begin{tikzpicture}[
        >=Latex,
        every node/.style={text=rtext},
        scale=1.0
    ]

        % Оси
        \draw[
            ->,
            draw=rtext,
            line width=0.8pt
        ]
        (0,0) -- (5.1,0)
        node[right] {$R$};

        \draw[
            ->,
            draw=rtext,
            line width=0.8pt
        ]
        (0,-2.0) -- (0,2.25)
        node[above] {$E$};

        % M < Mcr
        \draw[
            draw=rc3,
            line width=1.0pt,
            domain=0.55:4.7,
            samples=80,
            smooth,
            variable=\x
        ]
        plot ({\x},{1.25/\x});

        % M > Mcr
        \draw[
            draw=rtext,
            line width=0.9pt,
            domain=0.65:4.7,
            samples=80,
            smooth,
            variable=\x
        ]
        plot ({\x},{-1.05/\x});

        % M = Mcr
        \draw[
            dashed,
            draw=rtext,
            line width=0.7pt
        ]
        (0.45,0) -- (4.7,0);

        \node[
            right
        ]
        at (3.5,0.43)
        {$M<M_{\text{кр}}$};

        \node[
            right
        ]
        at (3.5,-0.38)
        {$M>M_{\text{кр}}$};

        \node[
            above
        ]
        at (2.2,0.03)
        {\small $M=M_{\text{кр}}$};

    \end{tikzpicture}

    \captionsetup{
        font=footnotesize,
        labelfont=bf,
        skip=5pt
    }

    \captionof{figure}{
        Качественная зависимость полной энергии
        ультрарелятивистской звезды от радиуса.
        В предельном случае $M=M_{\text{кр}}$
        коэффициент при $1/R$ обращается в нуль.
    }
    \label{fig:statphys_chandrasekhar_energy}
\end{rbox}


% ------------------------------------------------------------
% Давление ультрарелятивистского газа
% ------------------------------------------------------------

\subsection*{Давление ультрарелятивистского ферми-газа}

Теперь найдём давление более точно.

При $T=0$ энергия газа равна
%
\begin{equation*}
    E
    =
    2V
    \int_0^{p_F}
    pc\,
    \frac{d^3p}{(2\pi\hbar)^3}.
\end{equation*}

Используя
%
\begin{equation*}
    d^3p
    =
    4\pi p^2\,dp,
\end{equation*}
%
получаем
%
\begin{equation*}
    E
    =
    \frac{8\pi Vc}{(2\pi\hbar)^3}
    \int_0^{p_F}
    p^3\,dp.
\end{equation*}

Так как
%
\begin{equation*}
    \int_0^{p_F}
    p^3\,dp
    =
    \frac{p_F^4}{4},
\end{equation*}
%
имеем
%
\begin{equation*}
    E
    =
    \frac{Vc\,p_F^4}{4\pi^2\hbar^3}.
\end{equation*}

Плотность энергии
%
\begin{equation*}
    u
    =
    \frac{E}{V}
    =
    \frac{c\,p_F^4}{4\pi^2\hbar^3}.
\end{equation*}

Для ультрарелятивистского газа
%
\begin{equation*}
    P
    =
    \frac{u}{3}.
\end{equation*}

Следовательно,
%
\begin{equation*}
    \boxed{
        P
        =
        \frac{c\,p_F^4}
        {12\pi^2\hbar^3}
    }.
\end{equation*}

Используя
%
\begin{equation*}
    p_F
    =
    \hbar
    (3\pi^2n)^{1/3},
\end{equation*}
%
получаем
%
\begin{equation*}
    P
    =
    \frac{\hbar c}{4}
    (3\pi^2)^{1/3}
    n^{4/3}.
\end{equation*}

Итак,
%
\begin{equation*}
    \boxed{
        P
        =
        \frac{\hbar c}{4}
        (3\pi^2)^{1/3}
        n^{4/3}
    }.
\end{equation*}

По порядку величины
%
\begin{equation*}
    \boxed{
        P
        \sim
        \hbar c\,n^{4/3}
    }.
\end{equation*}


% ------------------------------------------------------------
% Электроны, протоны и нейтроны
% ------------------------------------------------------------

\subsection*{Соотношение концентраций электронов и нейтронов}

При больших плотностях становится существенной реакция
бета-равновесия
%
\begin{equation*}
    n
    \longleftrightarrow
    p
    +
    e^-.
\end{equation*}

Если нейтрино свободно покидают систему, их химический потенциал
можно положить равным нулю, поэтому условие химического равновесия:
%
\begin{equation*}
    \mu_n
    =
    \mu_p
    +
    \mu_e.
\end{equation*}

В предельно ультрарелятивистском приближении
%
\begin{equation*}
    \mu_i
    \simeq
    p_{F,i}c.
\end{equation*}

Следовательно,
%
\begin{equation*}
    p_{F,n}
    =
    p_{F,p}
    +
    p_{F,e}.
\end{equation*}

Из условия электрической нейтральности вещества
%
\begin{equation*}
    n_p
    =
    n_e.
\end{equation*}

Для частиц со спином $1/2$
%
\begin{equation*}
    n_i
    =
    \frac{p_{F,i}^3}
    {3\pi^2\hbar^3}.
\end{equation*}

Поэтому из
%
\begin{equation*}
    n_p
    =
    n_e
\end{equation*}
%
следует
%
\begin{equation*}
    p_{F,p}
    =
    p_{F,e}.
\end{equation*}

Тогда условие бета-равновесия принимает вид
%
\begin{equation*}
    p_{F,n}
    =
    2p_{F,e}.
\end{equation*}

Так как концентрация пропорциональна кубу импульса Ферми,
%
\begin{equation*}
    \frac{n_n}{n_e}
    =
    \left(
        \frac{p_{F,n}}
        {p_{F,e}}
    \right)^3
    =
    2^3.
\end{equation*}

Следовательно,
%
\begin{equation*}
    \boxed{
        n_n
        =
        8n_e
    }.
\end{equation*}

Или
%
\begin{equation*}
    \boxed{
        \frac{n_e}{n_n}
        =
        \frac{1}{8}
    }.
\end{equation*}


% ============================================================
% Итоги
% ============================================================

\section*{Основные результаты}

Для многоэлектронного атома в модели Томаса--Ферми
%
\begin{equation*}
    \boxed{
        R
        \sim
        a_0Z^{-1/3}
    }.
\end{equation*}

Для белого карлика с нерелятивистским вырожденным
электронным газом
%
\begin{equation*}
    \boxed{
        R
        \sim
        \frac{\hbar^2}
        {Gm_em_p^{5/3}}
        M^{-1/3}
    },
    \qquad
    R
    \propto
    M^{-1/3}.
\end{equation*}

Для ультрарелятивистского электронного газа возникает
характерная предельная масса
%
\begin{equation*}
    \boxed{
        M_{\text{кр}}
        \sim
        \frac{
            (\hbar c/G)^{3/2}
        }{
            m_p^2
        }
    }.
\end{equation*}

Давление полностью вырожденного ультрарелятивистского газа:
%
\begin{equation*}
    \boxed{
        P
        =
        \frac{\hbar c}{4}
        (3\pi^2)^{1/3}
        n^{4/3}
    }.
\end{equation*}

В предельно ультрарелятивистском случае при бета-равновесии
и электрической нейтральности
%
\begin{equation*}
    \boxed{
        n_n
        =
        8n_e
    }.
\end{equation*}


\end{document}